% 原创概念图；理想断连图的两维谱嵌入。
\begin{tikzpicture}[x=10mm,y=10mm,font=\small,text=BookInk,
  vertex/.style={circle,inner sep=0pt,minimum size=4pt,fill=BookTeal}]
 \node[vertex] (a) at (0,0) {};
 \node[vertex] (b) at (0.6,0.8) {};
 \node[vertex] (c) at (1.1,0) {};
 \node[vertex,rectangle,fill=BookBlue] (d) at (2.1,0.6) {};
 \node[vertex,rectangle,fill=BookBlue] (e) at (2.7,1.1) {};
 \node[vertex,rectangle,fill=BookBlue] (f) at (3.1,0.3) {};
 \draw[BookTeal,thick] (a) -- (b) -- (c) -- (a);
 \draw[BookBlue,thick] (d) -- (e) -- (f) -- (d);
 \draw[BookMuted,densely dashed] (c) -- (d);
 \node[below,align=center] at (1.5,-0.5) {两个分量\\虚线为可选弱连接};
 \draw[book arrow] (3.6,0.55) -- (5.1,0.55)
   node[midway,above,align=center] {谱嵌入\\行归一化};
 \draw[->,BookMuted] (5.7,-0.2) -- (8.8,-0.2) node[right] {$h_1$};
 \draw[->,BookMuted] (5.7,-0.2) -- (5.7,2.0) node[above] {$h_2$};
 \fill[BookTeal] (8,-0.2) circle[radius=3pt];
 \draw[BookBlue,fill=white,thick] (5.7,1.5) rectangle ++(0.10,0.10);
 \node[above,BookTeal] at (8,0) {分量一};
 \node[right,BookBlue] at (5.9,1.55) {分量二};
 \node[below,align=center] at (7.1,-0.5) {理想情形：\\每个分量汇于一个方向};
\end{tikzpicture}
